\subsection{Генерация и заполнение базы данных}

Код генерации таблиц базы данных был написан на языке SQL в PostgreSQL, исходный код приведен в приложении А. Таблицы первого уровня были заполнены также на языке SQL с помощью команды INSERT, исходный код приведен в приложении В. Код заполнения таблиц второго и третьего уровней был написан на языке Python, исходный код приведен в приложении С. Данные для таблиц были сгенерированы с помощью библиотеки faker из языка Python. 

Правила заполнения таблиц:
\begin{itemize}
\item 1 зоотехнику соответствует от 1 до 10 кормов;
\item в 1 рационе содержится от 10 до 15 записей;
\item 1 корове соответствует от 1 до 7 лактаций;
\item 1 корова питается по 1 - 30 рационам.
\end{itemize}

Количество записей в каждой из таблиц приведено в таблице \ref{tab:записи}. В сумме записей: 302 526.

\begin{table}[h]
\centering
\caption{Количество записей в таблицах БД}
\label{tab:записи}
\begin{tabular}{|l|c|}
\hline
\textbf{Название} & \textbf{Количество записей} \\
\hline
Breed & 10 \\
Day\_Times & 5 \\
Season & 4 \\
Lactation\_Phase & 5 \\
Feed\_Type & 17 \\
Livestock\_specialist & 8 \\
Cow & 6574 \\
Feeds & 40 \\
Diet & 38 \\
Lactation & 130910 \\
Feeding\_plan & 464 \\
Nutrition & 164451 \\
\hline
\end{tabular}
\end{table}